\documentclass[10pt,a4paper]{article}

% Encodage et langue
\usepackage[utf8]{inputenc}
\usepackage[T1]{fontenc}
\usepackage[french]{babel}
\usepackage{lmodern}
\usepackage{microtype}

% Mise en page
\usepackage{geometry}
\geometry{hmargin=2.2cm, vmargin=2.2cm}
\usepackage{setspace}
\setstretch{1.05}
\usepackage{fancyhdr}
\setlength{\headheight}{14pt}
\usepackage{titlesec}

% Tableaux et couleurs
\usepackage[table]{xcolor}
\usepackage{array}
\usepackage{booktabs}
\usepackage{tabularx}
\usepackage{enumitem}
\usepackage{hyperref}

\definecolor{SectionColor}{RGB}{0,76,153}
\definecolor{SoftBlue}{RGB}{238,244,252}
\definecolor{SoftGray}{RGB}{245,246,248}

\titleformat{\section}
  {\color{SectionColor}\normalfont\large\bfseries}{\thesection}{0.6em}{}
\titleformat{\subsection}
  {\color{SectionColor}\normalfont\normalsize\bfseries}{\thesubsection}{0.6em}{}

\setlist[itemize]{
  leftmargin=0.7cm,
  itemsep=0.15em,
  topsep=0.2em,
  parsep=0pt
}

\pagestyle{fancy}
\fancyhf{}
\lhead{\textbf{Maintenance Prédictive Industrielle}}
\rhead{\textbf{Version 1.0}}
\cfoot{\thepage}
\renewcommand{\headrulewidth}{0.4pt}

\newcolumntype{Y}{>{\raggedright\arraybackslash}X}
\newcommand{\important}[1]{\textbf{\color{SectionColor}#1}}

\begin{document}

% =========================
% Page de garde conservee
% =========================
\begin{titlepage}
\thispagestyle{empty}
\centering
\vspace*{2.5cm}

{\Huge\bfseries CAHIER DES CHARGES\par}
\vspace{0.8cm}
{\LARGE\bfseries Application de Maintenance Prédictive\par}
{\LARGE\bfseries Industrielle\par}

\vspace{1.2cm}
{\large Version 1.0\par}

\vspace{1.4cm}
\begin{flushleft}
\hspace{2.4cm}--- Étudiante: LAHDILI Hassna\\[0.18cm]
\hspace{2.4cm}--- École: Faculté Polydisciplinaire de Beni Mellal\\[0.18cm]
\hspace{2.4cm}--- Filière: Data Science et Sécurité des Systèmes d'Information\\[0.18cm]
\hspace{2.4cm}--- Tuteur entreprise: SERBOUTI Ibtissam\\[0.18cm]
\hspace{2.4cm}--- Tuteur académique: MAARIR Abdelkrim\\[0.18cm]
\hspace{2.4cm}--- Durée: 4 mois (du 02 mars 2026 au 29 juin 2026)\\[0.18cm]
\hspace{2.4cm}--- Objectif pédagogique: Développer une application de maintenance prédictive industrielle
\end{flushleft}

\vfill
\begin{flushleft}
\hspace{2.4cm}\textbf{Entreprise :} RAYNOV\\[0.25cm]
\hspace{2.4cm}\textbf{Date :} 20 avril 2026\\[0.25cm]
\hspace{2.4cm}\textbf{Statut :} À valider
\end{flushleft}
\end{titlepage}

\setcounter{page}{1}
\small

% =========================
% Page 1 du resume
% =========================
\section{Résumé exécutif}

\important{But de l'application.}
L'application réalisée est une plateforme web de maintenance prédictive industrielle.
Elle aide le responsable maintenance à surveiller l'état des machines, détecter les risques de panne, générer des alertes et suivre les interventions avant qu'un arrêt de production ne se produise.

\vspace{0.2cm}
\important{Problème traité.}
Dans une organisation industrielle, les pannes imprévues provoquent des arrêts coûteux.
La maintenance classique est souvent corrective ou préventive selon un calendrier fixe.
L'objectif du projet est de passer vers une maintenance pilotée par les données, les règles métier et l'intelligence artificielle.

\vspace{0.2cm}
\important{Objectifs principaux.}
\begin{itemize}
  \item Anticiper les pannes avec une cible d'au moins 7 jours d'avance.
  \item Réduire les arrêts non planifiés et améliorer la disponibilité des équipements.
  \item Centraliser les machines, capteurs, prédictions, alertes et interventions dans une même interface.
  \item Transformer l'expertise métier en règles exploitables par le modèle IA.
\end{itemize}

\subsection{Ce que le responsable peut comprendre rapidement}

\rowcolors{2}{SoftGray}{white}
\begin{tabularx}{\textwidth}{p{3.2cm}Y}
\toprule
\rowcolor{SoftBlue}
\textbf{Besoin} & \textbf{Réponse apportée par l'application} \\
\midrule
Superviser le parc & Dashboard avec KPIs, nombre de machines, alertes actives, machines à risque et état global du parc. \\
Détecter les risques & Analyse des données capteurs, application de règles métier et calcul de probabilité de panne. \\
Réagir aux alertes & Génération d'alertes par niveau de gravité, accusé de réception et historique de traitement. \\
Suivre les équipements & Gestion des machines, détails techniques, historique capteurs, prédictions et interventions. \\
Piloter l'IA & Page ML pour consulter les datasets, diagnostics, performances, drift, explications et retraining. \\
\bottomrule
\end{tabularx}

\newpage

% =========================
% Page 2 du resume
% =========================
\section{Fonctionnement général}

\subsection{Chaîne de fonctionnement}

\begin{enumerate}[leftmargin=0.8cm,itemsep=0.2em]
  \item L'utilisateur se connecte avec son compte sécurisé.
  \item Les machines industrielles sont créées et organisées par site, zone, ligne et composant.
  \item Les mesures capteurs sont enregistrées: température, vibration, pression, vitesse, courant, etc.
  \item Les règles métier sont évaluées pour repérer les situations anormales.
  \item Les résultats des règles deviennent des variables d'entrée pour le modèle IA.
  \item Le modèle calcule le risque de panne et le délai estimé avant défaillance.
  \item Les alertes sont créées, affichées, historisées et traitées par les utilisateurs autorisés.
\end{enumerate}

\subsection{Modules fonctionnels essentiels}

\rowcolors{2}{SoftGray}{white}
\begin{tabularx}{\textwidth}{p{3.5cm}Y}
\toprule
\rowcolor{SoftBlue}
\textbf{Module} & \textbf{Rôle dans l'application} \\
\midrule
Connexion & Protège l'accès avec authentification JWT et rôles utilisateurs. \\
Dashboard & Donne une vue temps réel du parc: KPIs, statuts, risques, alertes et évolution des probabilités de panne. \\
Machines & Permet de créer, modifier, consulter et organiser les équipements industriels. \\
Règles métier & Permet de définir des seuils et conditions de détection adaptés aux types de machines. \\
Alertes & Centralise les alertes prédictives, leur niveau de gravité, leur statut et leur accusé de réception. \\
Operations & Suit les actifs, interventions terrain, jobs d'ingestion et retours utilisés par la boucle MLOps. \\
ML & Présente les datasets, benchmarks, diagnostics, explications locales, drift et relance d'entraînement. \\
Utilisateurs & Gère les comptes et droits, réservé principalement à l'administrateur. \\
\bottomrule
\end{tabularx}

\subsection{Machines couvertes}

Le cahier des charges prévoit les machines industrielles suivantes: moteur, pompe, convoyeur, compresseur, ventilateur, générateur et broyeur.
L'application réalisée ajoute aussi le type \textit{four}, utile pour couvrir des cas industriels thermiques.

\newpage

% =========================
% Page 3 du resume
% =========================
\section{Intelligence métier et IA}

\subsection{Principe retenu}

L'application ne s'appuie pas uniquement sur un modèle IA.
Elle combine d'abord les règles métier issues de l'expertise terrain, puis utilise ces résultats pour enrichir la prédiction.
Cette approche rend le système plus compréhensible pour les responsables et plus utile pour les techniciens.

\rowcolors{2}{SoftGray}{white}
\begin{tabularx}{\textwidth}{p{3.4cm}Y}
\toprule
\rowcolor{SoftBlue}
\textbf{Étape} & \textbf{Explication} \\
\midrule
Règles métier & Exemple: température moteur trop élevée, vibration excessive, pression anormale, chute RPM. \\
Features IA & Les règles produisent un score de sévérité, un nombre de règles déclenchées, une marge au seuil et une confiance. \\
Prédiction & Le modèle estime la probabilité de panne, le niveau d'alerte et le risque machine. \\
Retour terrain & Les interventions et labels permettent d'améliorer progressivement le modèle. \\
\bottomrule
\end{tabularx}

\subsection{Données utilisées}

\begin{itemize}
  \item Données capteurs: température, vibrations, pression, vitesse de rotation, courant et timestamp.
  \item Données métier: type de machine, site, statut, règles actives, alertes et interventions.
  \item Données IA: probabilité de panne, niveau de risque, explication de prédiction et diagnostic modèle.
\end{itemize}

\subsection{Niveaux d'alerte}

\rowcolors{2}{SoftGray}{white}
\begin{tabularx}{\textwidth}{p{3cm}p{4cm}Y}
\toprule
\rowcolor{SoftBlue}
\textbf{Niveau} & \textbf{Seuil indicatif} & \textbf{Action attendue} \\
\midrule
Faible & Risque supérieur à 20\% & Surveillance. \\
Moyen & Risque supérieur à 40\% & Vérification planifiée. \\
Élevé & Risque supérieur à 60\% & Intervention à préparer. \\
Critique & Risque supérieur à 80\% & Action rapide et suivi prioritaire. \\
\bottomrule
\end{tabularx}

\section{Sécurité et droits}

L'accès est organisé par rôles: administrateur, technicien, expert et lecteur.
Les actions sensibles sont limitées: suppression de machine, gestion utilisateurs, modification des règles et consultation des logs.
Les mots de passe sont protégés, les sessions utilisent JWT, et la configuration production prévoit HTTPS, CORS restrictif et paramètres de sécurité renforcés.

\newpage

% =========================
% Page 4 du resume
% =========================
\section{Architecture, livrables et validation}

\subsection{Architecture technique}

\rowcolors{2}{SoftGray}{white}
\begin{tabularx}{\textwidth}{p{3.2cm}Y}
\toprule
\rowcolor{SoftBlue}
\textbf{Couche} & \textbf{Choix technique} \\
\midrule
Frontend & React avec Material UI pour une interface responsive et structurée. \\
Backend & FastAPI pour exposer les API REST, gérer l'authentification, les règles, les alertes et les prédictions. \\
Base de données & PostgreSQL pour stocker utilisateurs, machines, capteurs, alertes, prédictions, interventions et jobs. \\
Temps réel & WebSocket pour actualiser le dashboard lorsqu'un signal ou une intervention arrive. \\
ML/MLOps & Entraînement local, benchmarks datasets, diagnostics, drift, feedback terrain et retraining. \\
Déploiement & Docker Compose, fichiers d'environnement, migrations Alembic et configuration production HTTPS. \\
\bottomrule
\end{tabularx}

\subsection{Livrables principaux}

\begin{itemize}
  \item Application web complète: dashboard, machines, règles, alertes, opérations, ML et utilisateurs.
  \item API backend documentée avec endpoints d'authentification, supervision, prédiction, alertes et MLOps.
  \item Base PostgreSQL avec migrations et tables principales du domaine industriel.
  \item Pipeline ML avec modèle, entraînement, diagnostics, benchmarks et réentraînement à partir des retours terrain.
  \item Documentation projet: README, guide interface, plan de recette et configuration de déploiement.
\end{itemize}

\subsection{Critères de réussite}

\rowcolors{2}{SoftGray}{white}
\begin{tabularx}{\textwidth}{p{4.1cm}Y}
\toprule
\rowcolor{SoftBlue}
\textbf{Critère} & \textbf{Résultat attendu} \\
\midrule
Fonctionnel & Le dashboard, les machines, les règles, les alertes, l'authentification et les prédictions fonctionnent. \\
Performance & Dashboard sous 2 secondes et prédiction API sous 200 ms comme cible de référence. \\
IA & Rappel cible supérieur ou égal à 80\%, anticipation moyenne de 7 jours et suivi des faux positifs. \\
Sécurité & Accès protégé, rôles respectés, mots de passe hashés et configuration production sécurisée. \\
Recette & Tests automatisés et checklist manuelle pour valider conformité, interface, charge et qualité IA. \\
\bottomrule
\end{tabularx}

\subsection{Conclusion pour le responsable}

Cette application fournit un outil de pilotage maintenance clair: elle surveille les machines, détecte les anomalies, prédit les risques, déclenche les alertes et conserve l'historique des décisions.
Elle transforme les données industrielles et l'expertise métier en informations exploitables pour prioriser les interventions et réduire les arrêts imprévus.

\end{document}
